\documentclass[9pt,a4paper]{extarticle}
\usepackage[utf8]{inputenc}
\usepackage[T1]{fontenc}
\usepackage[brazil]{babel}
\usepackage[top=1.05cm, bottom=1.10cm, left=1.40cm, right=1.40cm]{geometry}
\usepackage{amsmath,amssymb}
\usepackage{booktabs}
\usepackage{array}
\usepackage{enumitem}
\usepackage{xcolor}
\usepackage{url}
\usepackage{hyperref}
\usepackage{titlesec}
\usepackage{microtype}

\definecolor{unbblue}{RGB}{0,51,102}
\definecolor{unbgreen}{RGB}{0,102,51}

\hypersetup{
    colorlinks=true,
    linkcolor=unbblue,
    urlcolor=unbblue,
    citecolor=unbblue
}

\pagestyle{plain}
\setlength{\parindent}{0pt}
\setlength{\parskip}{2.2pt}
\linespread{0.95}

\titleformat{\section}{\normalsize\bfseries\color{unbblue}}{\thesection.}{0.4em}{}[\titlerule]
\titlespacing*{\section}{0pt}{4.5pt}{2.5pt}
\titleformat{\subsection}{\small\bfseries\color{unbgreen}}{\thesubsection}{0.4em}{}
\titlespacing*{\subsection}{0pt}{3pt}{1pt}

\begin{document}

% Cabeçalho Institucional
\begin{center}
    {\large\bfseries\color{unbblue} UNIVERSIDADE DE BRAS\'{I}LIA (UnB) -- INSTITUTO DE CI\^{E}NCIAS EXATAS (IE)}\\
    {\normalsize\bfseries Departamento de Ci\^{e}ncia da Computa\c{c}\~{a}o (CIC) -- P\'{o}s-Gradua\c{c}\~{a}o em Inform\'{a}tica (PPGI)}\\
    {\footnotesize\bfseries Disciplina: PPGI0034 -- Redes Neurais e Aprendizado Profundo (2026/2) \quad|\quad Docente: Prof. Dr. D\'{i}bio Leandro Borges}\\
    {\footnotesize \textbf{Relat\'{o}rio Oficial de Especifica\c{c}\~{a}o e Auditoria do Uso de Ferramentas de Intelig\^{e}ncia Artificial (Projeto Individual 2)}}\\
    {\footnotesize \textbf{Autor:} Kelvin Rafael Paulino Buss (\texttt{kelvin.buss@aluno.unb.br}) \quad|\quad \textbf{GitHub:} \url{https://github.com/kelvinbuss/ProjetoFinalRNAP_UnB2026}}
\end{center}
\vspace{-5pt}
\noindent\textcolor{unbgreen}{\rule{\linewidth}{1.2pt}}
\vspace{-7pt}

\section{Escopo do Relat\'{o}rio e Especifica\c{c}\~{a}o dos Modelos de IA Utilizados}

Em cumprimento estrito ao edital do Projeto Individual 2 da disciplina PPGI0034 (\textit{``1) Listando e especificando como usou algum modelo de IA (descrever quais) no aux\'{i}lio de seu desenvolvimento. Seja espec\'{i}fico e detalhado. Duas (2) p\'{a}ginas''}), este documento detalha de forma transparente, audit\'{a}vel e exaustiva todas as ferramentas e modelos de Intelig\^{e}ncia Artificial utilizados ao longo do ciclo de vida da pesquisa intitulada \textbf{``Controle Din\^{a}mico de `Overthinking' em Modelos de Linguagem via Parada Precoce Calibrada por Entropia em Cadeias de Pensamento (\textit{Chain-of-Thought})''}, delimitando com rigor o que constitui concep\c{c}\~{a}o intelectual e realiza\c{c}\~{a}o autoral do mestrando e onde houve gera\c{c}\~{a}o assistida ou automa\c{c}\~{a}o computacional por IA.

O projeto envolveu duas categorias distintas de modelos de Intelig\^{e}ncia Artificial, com pap\'{e}is rigorosamente segregados:

\begin{enumerate}[leftmargin=1.4em, itemsep=1.5pt, topsep=1pt]
    \item \textbf{Modelos de Linguagem Objetos do Estudo Experimental (Executados Localmente via PyTorch/CUDA):}
    \begin{itemize}[leftmargin=1.2em, itemsep=1pt, topsep=0.5pt]
        \item \textbf{\texttt{DeepSeek-R1-Distill-Qwen-1.5B}} (DEEPSEEK-AI, 2025) e \textbf{\texttt{Qwen2.5-Math-1.5B-Instruct}} (YANG et al., 2024): \textit{Small Language Models} (SLMs) de 1,5 bilh\~{a}o de par\^{a}metros pertencentes ao estado da arte aberto para racioc\'{i}nio matem\'{a}tico via \textit{Chain-of-Thought} (CoT). Estes modelos \textbf{n\~{a}o} atuaram como assistentes de escrita, mas sim como o \textbf{objeto f\'{i}sico de investiga\c{c}\~{a}o cient\'{i}fica}, sendo submetidos a infer\^{e}ncia na GPU local do autor (NVIDIA GeForce RTX 3050 6\,GB VRAM) para coleta de probabilidades passo a passo, c\'{a}lculo da entropia de Shannon $H_t = -\sum_{v \in \mathcal{V}} p_\theta(v \mid x, y_{<t})\log p_\theta(v \mid x, y_{<t})$ e avalia\c{c}\~{a}o do decodificador proposto \textbf{LC-SC} nos benchmarks GSM8K e MATH-500.
    \end{itemize}
    \item \textbf{Ferramentas de IA Assistiva no Fluxo de Engenharia de Software e Editoria Acad\^{e}mica (\textit{Human-in-the-Loop}):}
    \begin{itemize}[leftmargin=1.2em, itemsep=1pt, topsep=0.5pt]
        \item \textbf{Ambiente Ag\^{e}ntico de Codifica\c{c}\~{a}o Assistida (Google Antigravity / Gemini Pro):} Utilizado sob supervis\~{a}o direta e cont\'{i}nua do autor como copiloto de engenharia de software e assistente de pesquisa bibliogr\'{a}fica para: (a) busca sistem\'{a}tica e download dos PDFs e reposit\'{o}rios GitHub oficiais do ICLR 2026; (b) estrutura\c{c}\~{a}o modular de scripts Python/PyTorch, tipagem est\'{a}tica e gera\c{c}\~{a}o de testes unit\'{a}rios automatizados (\texttt{pytest}); (c) codifica\c{c}\~{a}o de scripts de plotagem cient\'{i}fica vetorial (\texttt{matplotlib}/\texttt{seaborn} em 300 DPI); e (d) formata\c{c}\~{a}o tipogr\'{a}fica em LaTeX (\texttt{IEEEtran.cls}) com revis\~{a}o gramatical e padroniza\c{c}\~{a}o de refer\^{e}ncias ABNT \texttt{(AUTOR, ano)}.
    \end{itemize}
\end{enumerate}

\section{Delimita\c{c}\~{a}o entre Autoria Intelectual do Aluno e Apoio Operacional de IA}

A arquitetura metodol\'{o}gica do projeto seguiu o princ\'{i}pio de \textbf{Soberania Intelectual Humana}: todas as decis\~{o}es cient\'{i}ficas, formula\c{c}\~{o}es de hip\'{o}teses, escolhas de hiperpar\^{a}metros e interpreta\c{c}\~{o}es de resultados foram concebidas e validadas pelo autor, utilizando a IA assistiva estritamente como acelerador de implementa\c{c}\~{a}o e estrutura\c{c}\~{a}o. Especificamente, s\~{a}o de \textbf{autoria intelectual exclusiva do mestrando}:

\begin{itemize}[leftmargin=1.3em, itemsep=1.5pt, topsep=1pt]
    \item \textbf{Conex\~{a}o Epistemol\'{o}gica entre o Projeto 1, a Aula 11 e o ICLR 2026:} A decis\~{a}o de investigar a efici\^{e}ncia de racioc\'{i}nio em LLMs nasceu da experi\^{e}ncia pr\'{a}tica adquirida pelo autor no Projeto Individual 1 (onde observou a degenera\c{c}\~{a}o estoc\'{a}stica e la\c{c}os repetitivos em modelos autorregressivos \textit{nanoGPT}/\textit{Machado-GPT}) aliada ao estudo do paradigma \textit{Self-Consistency} apresentado pelo Prof. Dr. D\'{i}bio Leandro Borges na \textbf{Aula 11 (Slide 17)}.
    \item \textbf{Identifica\c{c}\~{a}o da Lacuna Cient\'{i}fica Cruzada (Gap nos Artigos do ICLR 2026):} Ao confrontar os artigos de (WU et al., 2026) e (HUANG et al., 2026), o autor identificou que ambos resolvem apenas metade do problema de custo $N \times L$: o \textit{Length-Filtered Vote} de (WU et al., 2026) controla o comprimento $L$ apenas \textit{post-hoc} (desperdi\c{c}ando 100\% dos tokens gerados em trilhas longas), enquanto o CaTS de (HUANG et al., 2026) controla apenas o n\'{u}mero de amostras $N$ e sofre colapso de calibra\c{c}\~{a}o em SLMs (queda de AUC de 64,31\% para 55,57\% na Tabela~1 do artigo no GSM8K) devido \`{a} infla\c{c}\~{a}o de probabilidades em loops de \textit{overthinking}.
    \item \textbf{Concep\c{c}\~{a}o do Algoritmo LC-SC (\textit{Length-Calibrated Self-Consistency}):} A ideia central de acoplar um gatilho de parada intratrilha guiado pela m\'{e}dia m\'{o}vel da entropia de Shannon ($\bar{H}_k(t) > \tau_H$ ap\'{o}s resposta candidata v\'{a}lida) com uma regra de vota\c{c}\~{a}o calibrada penalizada exponencialmente pelo excesso de comprimento $\exp(-\lambda \max(0, L_i - L^*)/\sigma_L)$ e parada intertrilhas ($S_{\max}(n) \ge \gamma$) constitui a contribui\c{c}\~{a}o algor\'{i}tmica original do autor.
    \item \textbf{Corre\c{c}\~{a}o de Infraestrutura Computacional (Li\c{c}\~{a}o Aprendida do Projeto 1):} No Projeto 1, a instala\c{c}\~{a}o padr\~{a}o do PyTorch operou em CPU (\texttt{+cpu}), exigindo 2h35min de treinamento. No Projeto 2, o autor determinou expressamente a configura\c{c}\~{a}o do ambiente virtual (\texttt{.venv}) com suporte nativo a \textbf{CUDA na GPU NVIDIA GeForce RTX 3050 (6\,GB VRAM)} e desenhou a estrat\'{e}gia de \textbf{Cache de Trilhas CoT em Parquet}, permitindo avaliar offline dezenas de combina\c{c}\~{o}es de $(\tau_H, \lambda, \gamma)$ sobre as distribui\c{c}\~{o}es reais do modelo em segundos, com reprodutibilidade exata.
\end{itemize}

\begin{table}[!htbp]
\caption{Matriz de Rastreabilidade e Divis\~{a}o de Responsabilidades: Autoria Humana vs. Aux\'{i}lio de IA}
\label{tab:matriz_ia}
\centering
\footnotesize
\renewcommand{\arraystretch}{1.04}
\begin{tabular}{p{2.8cm}p{5.8cm}p{5.6cm}p{2.5cm}}
\toprule
\textbf{Fase / Componente} & \textbf{Ideia, Realiza\c{c}\~{a}o e Decis\~{a}o do Autor (Aluno)} & \textbf{Apoio Operacional / Gera\c{c}\~{a}o da Ferramenta de IA} & \textbf{Mecanismo de Auditoria} \\
\midrule
\textbf{1. Revis\~{a}o ICLR 2026 e Lacuna de Pesquisa} & Sele\c{c}\~{a}o do tema de racioc\'{i}nio CoT (Aula 11, Slide 17); an\'{a}lise cr\'{i}tica da queda de AUC (55,57\%) no CaTS e do filtro \textit{post-hoc} de Wu et al. (2026). & Varredura nos anais do ICLR 2026; download automatizado dos 4 PDFs e dos 4 reposit\'{o}rios GitHub oficiais para a pasta \path{referencias/}. & Confer\^{e}ncia manual dos teoremas, tabelas e equa\c{c}\~{o}es nos PDFs originais. \\
\midrule
\textbf{2. Formula\c{c}\~{a}o do M\'{e}todo LC-SC ($N \times L$)} & Defini\c{c}\~{a}o dos 3 mecanismos ($L$-\textit{Stopping} por entropia, peso exponencial $L^*$ e $N$-\textit{Stopping}) e protocolo sem vazamento (\textit{cal}/\textit{test}). & Formaliza\c{c}\~{a}o da nota\c{c}\~{a}o matem\'{a}tica em LaTeX e estrutura\c{c}\~{a}o visual do pseudoc\'{o}digo (Algoritmo 1). & Verifica\c{c}\~{a}o anal\'{i}tica dos limites assint\'{o}ticos ($\lambda \to 0, \tau_H \to \infty$). \\
\midrule
\textbf{3. Engenharia Python, CUDA e Testes} & Diretriz de uso obrigat\'{o}rio da GPU RTX 3050 (CUDA FP16) e arquitetura de cache de trilhas passo a passo com entropia de Shannon. & Escrita modular de \path{src/decoders/}, \path{src/metrics/}, scripts de coleta/avalia\c{c}\~{a}o e su\'{i}te de testes \texttt{pytest}. & Execu\c{c}\~{a}o de testes unit\'{a}rios e inspe\c{c}\~{a}o de tensores e logs de GPU. \\
\midrule
\textbf{4. Experimentos, Figuras e Artigo IEEE} & Interpreta\c{c}\~{a}o da curva em U invertido, fronteira de Pareto e estudo de abla\c{c}\~{a}o; revis\~{a}o integral do texto cient\'{i}fico. & Gera\c{c}\~{a}o de scripts Matplotlib (300 DPI) e estrutura\c{c}\~{a}o tipogr\'{a}fica em \texttt{IEEEtran.cls} com cita\c{c}\~{o}es ABNT \texttt{(AUTOR, ano)}. & Compila\c{c}\~{a}o via \texttt{tectonic.exe} e checagem cruzada com os CSVs/JSONs. \\
\bottomrule
\end{tabular}
\end{table}

\section{Registro Detalhado de Prompts e Intera\c{c}\~{o}es no Desenvolvimento}

Para assegurar rastreabilidade metodol\'{o}gica completa (documentada integralmente em \path{docs/01_Relatorio_Evolucao_e_Prompts.md}), listam-se abaixo os principais comandos e diretrizes (\textit{prompts}) elaborados pelo autor para conduzir o assistente ag\^{e}ntico, seguidos do respectivo processamento t\'{e}cnico e valida\c{c}\~{a}o humana:

\begin{enumerate}[leftmargin=1.4em, itemsep=2pt, topsep=1pt]
    \item \textbf{Prompt 1 -- Mapeamento Estrat\'{e}gico do Edital, Slides das Aulas 10--12 e Artigos do ICLR 2026:}
    \begin{quote}
    \textit{``Vamos iniciar o Projeto 2. Veja todos os arquivos do Projeto 1 para aprender sobre os acertos e erros que fizemos e fazer melhor no Projeto 2. Na pasta do Projeto 2 já tem os slides das Aulas 10, 11 e 12 e as instruções do Projeto 2.''}
    \end{quote}
    \textbf{A\c{c}\~{a}o da IA e Curadoria do Autor:} O agente inspecionou o acervo do Projeto 1, extraiu os t\'{o}picos das Aulas 10, 11 e 12 do Prof. D\'{i}bio e mapeou 5 op\c{c}\~{o}es de pesquisa cruzadas com 20 artigos reais do ICLR 2026 (\path{opcoes_projeto_2_iclr2026.txt}). O autor analisou as alternativas e selecionou a \textbf{Op\c{c}\~{a}o 2 (Efici\^{e}ncia de Racioc\'{i}nio CoT e \textit{Test-Time Compute} em LLMs)}, por permitir conex\~{a}o direta com o Slide 17 da Aula 11 (\textit{Self-Consistency}), continuidade natural com a modelagem autorregressiva do Projeto 1 e execu\c{c}\~{a}o controlada na GPU RTX 3050 de 6\,GB.

    \item \textbf{Prompt 2 -- Consolida\c{c}\~{a}o da Etapa 1 (Proposta de 1 P\'{a}gina) e Download Integral das Refer\^{e}ncias ICLR 2026:}
    \begin{quote}
    \textit{``A Opção 2 é a melhor! Faça o documento de entrega da Etapa 1 (04/11/2026), deixe tudo muito bem organizado e já baixe na pasta de referências os 4 artigos do ICLR 2026 da Opção 2 e os códigos deles se tiverem disponíveis. Me explique passo a passo para eu entender todo o projeto, conseguir corrigir e auditar, além de ter uma expectativa de tempo para cada fase. Use subagentes em paralelo quando possível.''}
    \end{quote}
    \textbf{A\c{c}\~{a}o da IA e Curadoria do Autor:} Foram baixados localmente os 4 PDFs completos do ICLR 2026 --- (WU et al., 2026), (HUANG et al., 2026), (HAO et al., 2026) e (FENG et al., 2026) --- em \path{referencias/artigos_iclr2026_pdfs/} e clonados seus 4 reposit\'{o}rios oficiais do GitHub (\path{01_PKU-ML_CoT-Length}, \path{02_Chengsong-Huang_Self-Calibration_CaTS}, \path{03_tsinghua-fib-lab_RL-LLM-Reasoning} e \path{04_ByteDance-Seed_In-Place-TTT}) em \path{referencias/repos_iclr2026/}. A partir da inspe\c{c}\~{a}o direta dos c\'{o}digos e PDFs, o autor validou a formula\c{c}\~{a}o da proposta de 1 p\'{a}gina (\path{entrega_etapa1_04nov/}) e o plano de execu\c{c}\~{a}o paralela (\path{GUIA_PASSO_A_PASSO_PROJETO_2.md}).

    \item \textbf{Prompt 3 -- Execu\c{c}\~{a}o Completa do Pipeline Experimental em GPU, Artigo IEEE/ABNT e Apresenta\c{c}\~{a}o:}
    \begin{quote}
    \textit{``Pode continuar e executar todo o projeto seguindo o plano paralelo com subagentes! Configure o ambiente com CUDA na RTX 3050, implemente os módulos em src/ e os testes em tests/, rode a coleta de trilhas CoT e os benchmarks reais, gere todas as figuras de alta resolução, escreva o artigo IEEE completo (6 a 8 páginas com citações ABNT), o relatório de uso de IA de 2 páginas e os 4 slides da apresentação.''}
    \end{quote}
    \textbf{A\c{c}\~{a}o da IA e Curadoria do Autor:} O pipeline foi subdividido em frentes paralelas: (i) configura\c{c}\~{a}o do ambiente PyTorch com CUDA na RTX 3050 e download dos datasets GSM8K e MATH-500; (ii) implementa\c{c}\~{a}o do extrator de entropia token a token, dos baselines (\textit{Greedy CoT}, \textit{Standard SC}, \textit{Wu LFV} e \textit{Huang CaTS}) e do decodificador \textbf{LC-SC} validados com \texttt{pytest}; (iii) coleta de trilhas CoT na GPU, avalia\c{c}\~{a}o estat\'{i}stica com IC 95\% via \textit{bootstrap} e gera\c{c}\~{a}o das 4 figuras vetoriais; e (iv) reda\c{c}\~{a}o e compila\c{c}\~{a}o dos entreg\'{a}veis finais em LaTeX.
\end{enumerate}

\section{Protocolo de Verifica\c{c}\~{a}o Cient\'{i}fica, Preven\c{c}\~{a}o de Alucina\c{c}\~{o}es e Conformidade \'{E}tica}

O uso de ferramentas generativas na pesquisa acad\^{e}mica exige salvaguardas rigorosas contra dois riscos cl\'{a}ssicos: (i) cita\c{c}\~{o}es bibliogr\'{a}ficas inexistentes ou distorcidas e (ii) fabrica\c{c}\~{a}o ou vazamento de dados experimentais (\textit{data leakage}). Para eliminar ambos os riscos, o autor implementou um protocolo de auditoria em quatro camadas:

\begin{enumerate}[leftmargin=1.4em, itemsep=1.5pt, topsep=1pt]
    \item \textbf{Verifica\c{c}\~{a}o F\'{i}sica de Fontes Prim\'{a}rias do ICLR 2026:} Nenhuma refer\^{e}ncia bibliogr\'{a}fica foi citada sem que seu arquivo PDF integral estivesse fisicamente armazenado e verificado na pasta \path{referencias/artigos_iclr2026_pdfs/}. Todos os n\'{u}meros citados da literatura --- como o Teorema 4.3 e o Algoritmo 1 (p.~27) de (WU et al., 2026), a queda de AUC de 64,31\% para 55,57\% na Tabela 1 (p.~5) de (HUANG et al., 2026), e a an\'{a}lise de \textit{over-correction} na Figura 7 (Ap\^{e}ndice B.2, p.~19) de (HAO et al., 2026) --- foram auditados diretamente nas p\'{a}ginas correspondentes dos artigos originais.
    \item \textbf{Valida\c{c}\~{a}o de C\'{o}digo por Testes Unit\'{a}rios Determin\'{i}sticos (\texttt{pytest}):} Todo o c\'{o}digo gerado com aux\'{i}lio de IA em \path{src/} foi submetido a 31 testes automatizados distribu\'{i}dos em \path{tests/test_decoders_and_metrics.py}, \path{tests/test_data_and_traces_integrity.py}, \path{tests/test_benchmark_and_scientific_invariants.py} e \path{tests/test_deliverables_and_cross_doc_consistency.py}, verificando: (a) n\~{a}o-negatividade e limites te\'{o}ricos da entropia de Shannon $0 \le H_t \le \ln|\mathcal{V}|$; (b) redu\c{c}\~{a}o exata do peso $w_i$ ao valor $\hat{c}_i$ quando $L_i \le L^*$ e decaimento monot\^{o}nico para $L_i > L^*$; (c) acionamento correto do gatilho de $L$-\textit{Stopping} apenas quando $t > L_{\min} \wedge \bar{H}_k(t) > \tau_H \wedge \mathcal{D}_{\text{ans}} = 1$; e (d) equival\^{e}ncia num\'{e}rica exata do simulador sobre cache de trilhas em rela\c{c}\~{a}o \`{a} decodifica\c{c}\~{a}o online passo a passo.
    \item \textbf{Integridade Estat\'{i}stica e Separa\c{c}\~{a}o Estrita Calibra\c{c}\~{a}o/Teste:} Para impedir sobreajuste (\textit{overfitting}) dos hiperpar\^{a}metros $(\tau_H, L^*, \sigma_L, \lambda, \gamma)$ e do regressor isot\^{o}nico $g_{\text{cal}}$, o pipeline experimental isolou o subconjunto de calibra\c{c}\~{a}o $\mathcal{D}_{\text{cal}}$ do conjunto de teste $\mathcal{D}_{\text{test}}$. Todas as tabelas e figuras reportadas na Se\c{c}\~{a}o~IV do artigo IEEE refletem exclusivamente os logs reais de infer\^{e}ncia gravados em \path{results/}.
    \item \textbf{Responsabilidade Autoral Integral:} O autor declara que revisou, compreendeu e validou cada linha de c\'{o}digo, dedu\c{c}\~{a}o matem\'{a}tica e par\'{a}grafo textual que comp\~{o}em este trabalho, assumindo responsabilidade integral pela originalidade e exatid\~{a}o cient\'{i}fica dos resultados apresentados perante a banca examinadora da disciplina PPGI0034.
\end{enumerate}

\vspace{1pt}
\noindent\textcolor{unbblue}{\rule{\linewidth}{0.8pt}}
\begin{center}
\vspace{-4pt}
\scriptsize \textit{Documento gerado e compilado via \texttt{tectonic} em conformidade com as normas do Projeto Individual 2 -- PPGI0034 (Redes Neurais e Aprendizado Profundo, UnB, 2026/2).}
\end{center}

\end{document}
